% Workflow di progetto: dallo schematico all'ordine JLCPCB.
% Derivato da WORKFLOW.md: le modifiche vanno fatte prima li', poi qui.
%
% PDF:  firmware/MCU/produce-pdf-from-tex.sh WORKFLOW   ->  WORKFLOW.pdf
\documentclass[11pt,a4paper]{article}

\usepackage[T1]{fontenc}
\usepackage[utf8]{inputenc}
\usepackage{lmodern}
\usepackage[italian]{babel}
\usepackage{geometry}
\usepackage{listings}
\usepackage{xcolor}
\usepackage{booktabs}
\usepackage{tabularx}
\usepackage{array}
\usepackage{enumitem}
\usepackage{fancyhdr}
\usepackage{titlesec}
\usepackage{parskip}
\usepackage{tikz}
\usetikzlibrary{positioning,arrows.meta}
\usepackage[hidelinks]{hyperref}

\newcommand{\projname}{VERA-MODULE-RBL-XE}

\geometry{top=25mm, bottom=25mm, left=25mm, right=25mm}

% ---------------------------------------------------------------------------
% Colori e stili
% ---------------------------------------------------------------------------
\definecolor{codebg}{RGB}{245,245,245}
\definecolor{codeframe}{RGB}{200,200,200}
\definecolor{stepfill}{RGB}{232,240,252}
\definecolor{stepline}{RGB}{60,90,150}
\definecolor{stopfill}{RGB}{252,232,232}
\definecolor{okfill}{RGB}{232,248,236}

\lstdefinestyle{shell}{
    backgroundcolor=\color{codebg},
    frame=single,
    framerule=0.4pt,
    rulecolor=\color{codeframe},
    basicstyle=\ttfamily\small,
    breaklines=true,
    showstringspaces=false,
    xleftmargin=6pt,
    xrightmargin=6pt,
    aboveskip=8pt,
    belowskip=8pt,
}

% Nomi di file: a capo dopo i trattini, i nomi qui sono lunghi
% (solo dopo i trattini veri: niente sillabazione dentro i nomi)
\makeatletter
\newcommand{\file}[1]{{\ttfamily\hyphenchar\font=`\-\relax\language=\l@nohyphenation #1}}
\makeatother
\setlength{\emergencystretch}{3em}
\newcolumntype{L}{>{\raggedright\arraybackslash}X}

\pagestyle{fancy}
\fancyhf{}
\lhead{\small \projname}
\rhead{\small Workflow di progetto}
\cfoot{\small \thepage}
\renewcommand{\headrulewidth}{0.4pt}

\titleformat{\section}{\Large\bfseries\color{stepline}}{\thesection.}{0.6em}{}
\titleformat{\subsection}{\large\bfseries}{\thesubsection}{0.6em}{}

\title{\textbf{Workflow di progetto}\\[4pt]
       \large Dallo schematico all'ordine JLCPCB\\[2pt]
       \normalsize \projname}
\author{RetroBitLab}
\date{\today}

\begin{document}

\maketitle
\thispagestyle{empty}

\noindent
Strumenti: KiCad~5 (Eeschema, CvPcb, Pcbnew) per schematico e PCB, la libreria
\nolinkurl{~/Progetti/RetroBitLab-Library} per simboli, footprint e
modelli 3D, e gli script del repository per tutto ciò che si produce (PDF,
Gerber, 3D, BOM, CPL).

In questa fase \textbf{non} si esportano a mano PDF, BOM, Gerber o CPL da
KiCad: li genera tutti \file{output\_job.sh} al passo~10, sempre dai sorgenti.

\begin{figure}[h]
\centering
\begin{tikzpicture}[
    node distance=4mm and 8mm,
    wfstep/.style={draw=stepline, fill=stepfill, rounded corners=2pt,
                 minimum width=50mm, minimum height=7mm, font=\small, align=center},
    check/.style={wfstep, fill=stopfill},
    done/.style={wfstep, fill=okfill},
    arr/.style={-{Stealth[length=2mm]}, stepline, thick},
]
\node[wfstep]  (s1)  {1. Schematico e annotazione};
\node[check, below=of s1] (s2) {2. Disponibilità dei componenti};
\node[wfstep, right=16mm of s2]  (lib) {Aggiornare la libreria};
\node[check, below=of s2] (s3) {3. ERC};
\node[wfstep,  below=of s3] (s4) {4. Assegnamento dei footprint};
\node[wfstep,  below=of s4] (s5) {5. Salvataggio della netlist};
\node[wfstep,  below=of s5] (s6) {6. PCB: netlist e netclass};
\node[wfstep,  below=of s6] (s7) {7. Sbroglio};
\node[check, below=of s7] (s8) {8. Riempimento aree e DRC};
\node[wfstep,  below=of s8] (s9) {9. Salvataggio del PCB};
\node[wfstep,  below=of s9] (s10) {10. \texttt{./output\_job.sh <versione>}};
\node[done,  below=of s10] (s11) {11. Consegna: zip, BOM, CPL};

\foreach \a/\b in {s1/s2, s2/s3, s3/s4, s4/s5, s5/s6, s6/s7, s7/s8, s8/s9, s9/s10, s10/s11}
    \draw[arr] (\a) -- (\b);
\draw[arr] (s2) -- node[below, font=\scriptsize] {no} (lib);
\draw[arr] (lib) |- (s1);
\draw[arr] (s8.west) -- ++(-6mm,0) |- node[pos=0.25, left, font=\scriptsize, align=right] {errori\\o warning} (s7.west);
\end{tikzpicture}
\caption{Sequenza dei passi. In rosso i controlli che possono rimandare indietro.}
\end{figure}

\clearpage
\tableofcontents
\clearpage

% ---------------------------------------------------------------------------
\section{Schematico e annotazione}

\begin{itemize}[leftmargin=*]
  \item Si disegna lo schematico in Eeschema (KiCad~5): foglio principale
        \file{VERA-MODULE-RBL.sch} e fogli gerarchici.
  \item Ogni componente da montare ha il campo \textbf{\file{LCSC}} con il
        codice del pezzo (es.\ \file{C17713}). La BOM raggruppa per valore,
        footprint e LCSC; un componente senza LCSC non viene montato da
        JLCPCB e viene tolto anche dalla CPL.
  \item Si annota tutto (\emph{Strumenti $\rightarrow$ Annota schematico}).
        Nessun riferimento deve restare con \file{?}: \file{sch2bomxml.py} si
        ferma se ne trova uno, o se due parti diverse hanno lo stesso
        riferimento.
\end{itemize}

% ---------------------------------------------------------------------------
\section{Disponibilità dei componenti}

Subito dopo l'annotazione, prima di andare avanti:

\begin{lstlisting}[style=shell]
./check_jlcpcb_stock.sh          # 5 schede
./check_jlcpcb_stock.sh -n 10    # 10 schede
\end{lstlisting}

Lo script costruisce la BOM dallo schematico in una cartella temporanea (non
scrive file del progetto) e la confronta con il catalogo JLCPCB in tempo reale.
Per ogni riga riporta stock, quantità necessaria (con scarto e minimo), tipo
Basic / Preferred / Extended e costo.

\begin{center}
\begin{tabularx}{\linewidth}{>{\ttfamily}l L c}
\toprule
\normalfont\textbf{Stato} & \textbf{Significato} & \textbf{Blocca} \\
\midrule
OK      & stock sufficiente & no \\
LOW     & sufficiente ora, ma meno di 5 volte il necessario & no (avviso) \\
GLOBAL  & stock insufficiente, ma il codice è in \file{jlcpcb-global-sourcing.txt} & no \\
SHORT   & stock minore del necessario & sì \\
OUT     & stock zero & sì \\
MISSING & codice non presente nel catalogo JLCPCB & sì \\
\bottomrule
\end{tabularx}
\end{center}

Per le righe che bloccano propone alternative con stock sufficiente (stessa
categoria e package, stesso valore per i passivi, Basic prima). Esce con~0 se
tutto si può montare, 1 se qualcosa va cambiato, 2 se il servizio non risponde.

Se un componente è attivo ma JLCPCB ne ha pochi, lo si può far comprare a
JLCPCB con il \textbf{Global Sourcing}: si aggiunge il codice LCSC a
\file{jlcpcb-global-sourcing.txt} (uno per riga, \file{\#} per il commento) e
all'ordine si sceglie Global Sourcing per quella riga.

\subsection*{Aggiornare la libreria, se necessario}

Quando un componente va sostituito:

\begin{enumerate}[leftmargin=*]
  \item in \file{RetroBitLab-Library} si crea o si aggiorna il simbolo
        (\file{.lib}) con footprint, datasheet e campo \file{LCSC}, e la sua
        documentazione (\file{.dcm}); per i passivi si usa la libreria della
        famiglia giusta (\file{CCAP0805}, \file{CRES0805}, \dots). Se cambia il
        package servono anche footprint e modello 3D;
  \item si verifica sul datasheet piedinatura e land pattern rispetto al
        footprint;
  \item si fa il commit della libreria;
  \item nello schematico si sostituisce il simbolo e si torna al passo~2
        finché il controllo non passa.
\end{enumerate}

% ---------------------------------------------------------------------------
\section{ERC}

\emph{Ispeziona $\rightarrow$ Controllo regole elettriche}. Si correggono tutti
gli errori e i warning prima di proseguire.

% ---------------------------------------------------------------------------
\section{Assegnamento dei footprint}

Con CvPcb ogni simbolo deve avere il suo footprint. Per le parti JLCPCB il
footprint deve corrispondere al package del codice LCSC.

% ---------------------------------------------------------------------------
\section{Salvataggio della netlist}

Si genera e si salva la netlist \file{VERA-MODULE-RBL.net} (\emph{Strumenti
$\rightarrow$ Genera netlist}, formato Pcbnew). \textbf{Nessun PDF e nessuna
BOM} in questa fase.

% ---------------------------------------------------------------------------
\section{PCB: netlist e netclass}

\begin{itemize}[leftmargin=*]
  \item In Pcbnew (KiCad~5) si importa la netlist (\emph{Strumenti
        $\rightarrow$ Carica netlist}).
  \item Le \textbf{netclass si definiscono e si assegnano nel PCB}
        (\emph{File $\rightarrow$ Impostazioni scheda $\rightarrow$ Classi di
        rete}): con KiCad~5 non esiste back annotation dal PCB allo
        schematico, quindi l'unico posto dove le netclass sono definite è il
        \file{.kicad\_pcb}. Dopo ogni nuovo import della netlist si controlla
        che le net nuove o rinominate siano nella netclass giusta.
\end{itemize}

% ---------------------------------------------------------------------------
\section{Sbroglio}

A mano, oppure con fastroute tramite \file{tools/route.sh} (vedi il commento in
testa allo script). Il risultato si importa come sessione Specctra.

% ---------------------------------------------------------------------------
\section{Riempimento delle aree e DRC}

\begin{enumerate}[leftmargin=*]
  \item Al termine dello sbroglio si riempiono tutte le aree (tasto \file{B}).
  \item Si lancia il DRC (\emph{Ispeziona $\rightarrow$ Controllo regole di
        progettazione}).
  \item Si correggono \textbf{tutti} gli errori e i warning, poi si ripetono i
        punti 1 e 2 finché il DRC è pulito.
\end{enumerate}

% ---------------------------------------------------------------------------
\section{Salvataggio del PCB}

Il \file{.kicad\_pcb} si salva sempre corretto (DRC pulito) e \textbf{con le
aree riempite}: Gerber, 3D e CPL vengono generati dal file così come è salvato.

\file{REFILL\_ZONES=1 ./output\_job.sh <versione>} riempie di nuovo le aree solo
in memoria, senza modificare il file: è un ripiego, non sostituisce il passo~8.

% ---------------------------------------------------------------------------
\section{Generazione dei file di produzione}

\begin{lstlisting}[style=shell]
./check_jlcpcb_stock.sh          # di nuovo, lo stock cambia nel tempo
./output_job.sh <versione>       # es. ./output_job.sh 4
\end{lstlisting}

\file{output\_job.sh}, in ordine:

\begin{enumerate}[leftmargin=*]
  \item XML della BOM dallo schematico (\file{sch2bomxml.py})
        $\rightarrow$ \file{VERA-MODULE-RBL.xml};
  \item BOM raggruppata (\file{generate\_bom.py})
        $\rightarrow$ \file{VERA-MODULE-RBL.csv};
  \item \textbf{BOM JLCPCB} (\file{jlcpcb-check-bom.py})
        $\rightarrow$ \file{<NOME>-v<versione>-JLCPCB-BOM.csv};
  \item \textbf{CPL JLCPCB} dal PCB (\file{run\_generate\_cpl.sh},
        \file{clean\_virtual.sh}) $\rightarrow$
        \file{<NOME>-v<versione>-JLCPCB-CPL.csv}, poi
        \file{jlcpcb-sync-cpl.py} toglie dalla CPL ciò che non è in BOM e si
        ferma se manca qualcosa;
  \item \textbf{Gerber e forature} con Pcbnew~5 e le impostazioni di stampa
        salvate nel PCB (\file{generate\_gerbers.py}) $\rightarrow$
        \file{production/} (svuotata prima);
  \item \textbf{PDF dello schematico} con eeplot $\rightarrow$
        \file{VERA-MODULE-RBL.pdf};
  \item \textbf{scheda in 3D}, STEP e VRML, con kicad-cli $\rightarrow$
        \file{images/VERA-MODULE-RBL.step} e \file{images/VERA-MODULE-RBL.wrl};
  \item zip del contenuto di \file{production/} $\rightarrow$
        \file{<NOME>-v<versione>.zip}.
\end{enumerate}

Si ferma al primo errore. \file{<NOME>} è \file{\projname} per questo progetto.

% ---------------------------------------------------------------------------
\section{Consegna}

Il progetto da mandare in produzione è composto da \textbf{3 file}:

\begin{center}
\begin{tabularx}{\linewidth}{>{\raggedright\arraybackslash\ttfamily\small}p{0.42\linewidth} L L}
\toprule
\normalfont\normalsize\textbf{File} & \textbf{Contenuto} & \textbf{Dove si carica su JLCPCB} \\
\midrule
<NOME>-v<versione>.zip             & Gerber e forature     & ordine PCB \\
<NOME>-v<versione>-JLCPCB-BOM.csv  & BOM                   & PCB Assembly $\rightarrow$ BOM \\
<NOME>-v<versione>-JLCPCB-CPL.csv  & posizioni e rotazioni & PCB Assembly $\rightarrow$ CPL \\
\bottomrule
\end{tabularx}
\end{center}

All'ordine:

\begin{itemize}[leftmargin=*]
  \item nell'anteprima di montaggio si controlla la rotazione dei componenti,
        soprattutto quelli nuovi o con footprint cambiato;
  \item per le righe elencate come \file{GLOBAL} dal controllo stock si sceglie
        Global Sourcing.
\end{itemize}

Zip, BOM e CPL versionati non vanno nel repository; ci vanno invece i sorgenti
(schematici, netlist, PCB), \file{VERA-MODULE-RBL.xml} e \file{.csv},
\file{production/}, il PDF e i modelli 3D.

\end{document}
